% !TeX encoding = UTF-8
% !TeX program = xelatex
% !TeX spellcheck = en_US

\documentclass[degree=master]{sysuthesis}
  % 学位 degree:
  %   doctor | master | bachelor | postdoc
  % 学位类型 degree-type:
  %   academic（默认）| professional


% 论文基本配置，加载宏包等全局配置
% !TeX root = ./main.tex

% 论文基本信息配置

\sysusetup{
  %******************************
  % 注意：
  %   1. 配置里面不要出现空行
  %   2. 不需要的配置信息可以删除
  %   3. 建议先阅读文档中所有关于选项的说明
  %******************************
  %
  % 输出格式
  %   选择打印版（print）或用于提交的电子版（electronic），前者会插入空白页以便直接双面打印
  %
  output = print,
  %
  % 标题
  %   可使用“\\”命令手动控制换行
  %
  title  = {中山大学学位论文 \LaTeX{} 模板\\Sun Yat-sen University Thesis \LaTeX{} Template},
  % title* = {An Introduction to \LaTeX{} Thesis Template of Tsinghua
  %           University v\version},
  %
  % 学位
  %   1. 学术型
  %      - 中文
  %        需注明所属的学科门类，例如：
  %        哲学、经济学、法学、教育学、文学、历史学、理学、工学、农学、医学、
  %        军事学、管理学、艺术学
  %      - 英文
  %        博士：Doctor of Philosophy
  %        硕士：
  %          哲学、文学、历史学、法学、教育学、艺术学门类，公共管理学科
  %          填写“Master of Arts“，其它填写“Master of Science”
  %   2. 专业型
  %      直接填写专业学位的名称，例如：
  %      教育博士、工程硕士等
  %      Doctor of Education, Master of Engineering
  %   3. 本科生不需要填写
  %
  degree-name  = {博士},
  %
  % 培养单位
  %   填写所属院系的全名
  %
  department = {计算机学院},
  %
  %
  discipline  = {计算机科学与技术},
  %
  % 姓名
  author  = {研究生},
  %
  % 指导教师
  %   中文姓名和职称之间以英文逗号“,”分开，下同
  supervisor  = {某某某, 教授},
  %
  % 联合指导教师
  % joint-supervisor  = {某某某, 教授},
  %
  % 答辩委员会成员
  committee = {老师甲\ 教授, 老师乙\ 副教授, 老师丙\ 研究员, 老师丁\ 副研究员},
  % 日期
  %   使用 ISO 格式；默认为当前时间
  %
  % date = {2019-07-07},
  %
  %
  % 密级和年限
  %   秘密, 机密, 绝密
  %
  % secret-level = {秘密},
  % secret-year  = {10},
  %
}

% 载入所需的宏包

% 可以使用 nomencl 生成符号和缩略语说明
% \usepackage{nomencl}
% \makenomenclature

% 表格加脚注
\usepackage{threeparttable}

% 表格中支持跨行
\usepackage{multirow}

% 固定宽度的表格。放在 hyperref 之前的话，tabularx 里的 footnote 显示不出来。
% \usepackage{tabularx}

% 跨页表格
% \usepackage{longtable}

% 量和单位
\usepackage{siunitx}

% 定理类环境宏包
\usepackage{amsthm}
% 也可以使用 ntheorem
% \usepackage[amsmath,thmmarks,hyperref]{ntheorem}

% 参考文献使用 BibTeX + natbib 宏包
% 顺序编码制
\usepackage[sort]{natbib}
\bibliographystyle{sysuthesis-numeric}

% 著者-出版年制
% \usepackage{natbib}
% \bibliographystyle{sysuthesis-author-year}

% 本科生参考文献的著录格式
% \usepackage[sort]{natbib}
% \bibliographystyle{sysuthesis-bachelor}

% 参考文献使用 BibLaTeX 宏包
% \usepackage[backend=biber,style=sysuthesis-numeric]{biblatex}
% \usepackage[backend=biber,style=sysuthesis-author-year]{biblatex}
% \usepackage[backend=biber,style=apa]{biblatex}
% \usepackage[backend=biber,style=mla-new]{biblatex}
% 声明 BibLaTeX 的数据库
% \addbibresource{ref/refs.bib}

% 定义所有的图片文件在 figures 子目录下
\graphicspath{{figures/}}

% 数学命令
\newcommand\dif{\mathop{}\!\mathrm{d}}  % 微分符号

% hyperref 宏包在最后调用
\usepackage{hyperref}



\begin{document}

% 封面
\maketitle

% % 学位论文指导小组、公开评阅人和答辩委员会名单
% % !TeX root = ../thuthesis-example.tex

\begin{committee}[name={学位论文指导小组、公开评阅人和答辩委员会名单}]

  \newcolumntype{C}[1]{@{}>{\centering\arraybackslash}p{#1}}

  \section*{指导小组名单}

  \begin{center}
    \begin{tabular}{C{3cm}C{3cm}C{9cm}@{}}
      李XX & 教授     & 清华大学 \\
      王XX & 副教授   & 清华大学 \\
      张XX & 助理教授 & 清华大学 \\
    \end{tabular}
  \end{center}


  \section*{公开评阅人名单}

  \begin{center}
    \begin{tabular}{C{3cm}C{3cm}C{9cm}@{}}
      刘XX & 教授   & 清华大学                    \\
      陈XX & 副教授 & XXXX大学                    \\
      杨XX & 研究员 & 中国XXXX科学院XXXXXXX研究所 \\
    \end{tabular}
  \end{center}


  \section*{答辩委员会名单}

  \begin{center}
    \begin{tabular}{C{2.75cm}C{2.98cm}C{4.63cm}C{4.63cm}@{}}
      主席 & 赵XX                  & 教授                    & 清华大学       \\
      委员 & 刘XX                  & 教授                    & 清华大学       \\
          & \multirow{2}{*}{杨XX} & \multirow{2}{*}{研究员} & 中国XXXX科学院 \\
          &                       &                         & XXXXXXX研究所  \\
          & 黄XX                  & 教授                    & XXXX大学       \\
          & 周XX                  & 副教授                  & XXXX大学       \\
      秘书 & 吴XX                  & 助理研究员              & 清华大学       \\
    \end{tabular}
  \end{center}

\end{committee}



% 也可以导入 Word 版转的 PDF 文件
% \begin{committee}[file=figures/committee.pdf]
% \end{committee}


% 声明

% 生成的声明页是否要插入页眉和页脚（默认 empty）
% 仅在需要进行电子签名时，才需要打开这一选项
% 插入的扫描声明页总是会生成页眉（研究生）和页脚，不受这一选项影响
% \statement[page-style=plain]
% 将签字扫描后的声明文件 scan-statement.pdf 替换原始页面
% \statement[file=scan-statement.pdf]


% % 使用授权的说明
\statementpage
% % 将签字扫描后授权文件 scan-copyright.pdf 替换原始页面
% % \copyrightpage[file=scan-copyright.pdf]

\frontmatter
% !TeX root = ../thuthesis-example.tex

% 中英文摘要和关键字

\begin{abstract}
  论文的摘要是对论文研究内容和成果的高度概括。
  摘要应对论文所研究的问题及其研究目的进行描述，对研究方法和过程进行简单介绍，对研究成果和所得结论进行概括。
  摘要应具有独立性和自明性，其内容应包含与论文全文同等量的主要信息。
  使读者即使不阅读全文，通过摘要就能了解论文的总体内容和主要成果。

  论文摘要的书写应力求精确、简明。
  切忌写成对论文书写内容进行提要的形式，尤其要避免“第 1 章……；第 2 章……；……”这种或类似的陈述方式。

  关键词是为了文献标引工作、用以表示全文主要内容信息的单词或术语。
  关键词不超过 5 个，每个关键词中间用分号分隔。

  % 关键词用“英文逗号”分隔，输出时会自动处理为正确的分隔符
  \sysusetup{
    keywords = {关键词 1, 关键词 2, 关键词 3, 关键词 4, 关键词 5},
  }
\end{abstract}

\begin{abstract*}
  An abstract of a dissertation is a summary and extraction of research work and contributions.
  Included in an abstract should be description of research topic and research objective, brief introduction to methodology and research process, and summarization of conclusion and contributions of the research.
  An abstract should be characterized by independence and clarity and carry identical information with the dissertation.
  It should be such that the general idea and major contributions of the dissertation are conveyed without reading the dissertation.

  An abstract should be concise and to the point.
  It is a misunderstanding to make an abstract an outline of the dissertation and words “the first chapter”, “the second chapter” and the like should be avoided in the abstract.

  Keywords are terms used in a dissertation for indexing, reflecting core information of the dissertation.
  An abstract may contain a maximum of 5 keywords, with semi-colons used in between to separate one another.

  % Use comma as seperator when inputting
  \sysusetup{
    keywords* = {keyword 1, keyword 2, keyword 3, keyword 4, keyword 5},
  }
\end{abstract*}


% 目录
\tableofcontents

% 插图和附表清单
\listoffiguresandtables  % 插图和附表清单
% \listoffigures           % 插图清单
% \listoftables            % 附表清单

% 符号对照表
% !TeX root = ../main.tex

\begin{denotation}[3cm]
  \item[] {\fangsong (请按照字母顺序(A-Z)排序)}
  \item[$\kappa$] 传输系数 (Transmission Coefficient)
  \item[$\nu_i$] 虚频 (Imaginary Frequency)
  \item[\textit{ab initio}] 基于第一原理的量子化学计算方法，常称从头算法
  \item[As-PPT] 聚苯基不对称三嗪
  \item[DFT] 密度泛函理论 (Density Functional Theory)
  \item[$E_a$] 化学反应的活化能 (Activation Energy)
  \item[$\increment G^\neq$] 活化自由能(Activation Free Energy)
  \item[HPCE] 高效毛细管电泳色谱 (High Performance Capillary lectrophoresis)
  \item[HPLC] 高效液相色谱 (High Performance Liquid Chromatography)
  \item[IRC] 内禀反应坐标 (Intrinsic Reaction Coordinates)
  \item[LC-MS] 液相色谱-质谱联用 (Liquid chromatography-Mass Spectrum)
  \item[MPBI] 聚苯并咪唑模型化合物，N-苯基苯并咪唑
  \item[MPI] 聚酰亚胺模型化合物，N-苯基邻苯酰亚胺
  \item[MSPPT] 聚苯基对称三嗪模型化合物，2,4,6-三苯基-1,3,5-三嗪
  \item[ONIOM] 分层算法 (Our own N-layered Integrated molecular Orbital and molecular Mechanics)
  \item[PBI] 聚苯并咪唑
  \item[PDT] 热分解温度
  \item[PES] 势能面 (Potential Energy Surface)
  \item[PI] 聚酰亚胺
  \item[PPQ] 聚苯基喹噁啉
  \item[PY] 聚吡咙
  \item[PMDA-BDA] 均苯四酸二酐与联苯四胺合成的聚吡咙薄膜
  \item[SCF] 自洽场 (Self-Consistent Field)
  \item[S-PPT] 聚苯基对称三嗪
  \item[TIC] 总离子浓度 (Total Ion Content)
  \item[TS] 过渡态 (Transition State)
  \item[TST] 过渡态理论 (Transition State Theory)
  \item[ZPE] 零点振动能 (Zero Vibration Energy)
\end{denotation}



% 也可以使用 nomencl 宏包，需要在导言区
% \usepackage{nomencl}
% \makenomenclature

% 在这里输出符号说明
% \printnomenclature[3cm]

% 在正文中的任意为都可以标题
% \nomenclature{PI}{聚酰亚胺}
% \nomenclature{MPI}{聚酰亚胺模型化合物，N-苯基邻苯酰亚胺}
% \nomenclature{PBI}{聚苯并咪唑}
% \nomenclature{MPBI}{聚苯并咪唑模型化合物，N-苯基苯并咪唑}
% \nomenclature{PY}{聚吡咙}
% \nomenclature{PMDA-BDA}{均苯四酸二酐与联苯四胺合成的聚吡咙薄膜}
% \nomenclature{MPY}{聚吡咙模型化合物}
% \nomenclature{As-PPT}{聚苯基不对称三嗪}
% \nomenclature{MAsPPT}{聚苯基不对称三嗪单模型化合物，3,5,6-三苯基-1,2,4-三嗪}
% \nomenclature{DMAsPPT}{聚苯基不对称三嗪双模型化合物（水解实验模型化合物）}
% \nomenclature{S-PPT}{聚苯基对称三嗪}
% \nomenclature{MSPPT}{聚苯基对称三嗪模型化合物，2,4,6-三苯基-1,3,5-三嗪}
% \nomenclature{PPQ}{聚苯基喹噁啉}
% \nomenclature{MPPQ}{聚苯基喹噁啉模型化合物，3,4-二苯基苯并二嗪}
% \nomenclature{HMPI}{聚酰亚胺模型化合物的质子化产物}
% \nomenclature{HMPY}{聚吡咙模型化合物的质子化产物}
% \nomenclature{HMPBI}{聚苯并咪唑模型化合物的质子化产物}
% \nomenclature{HMAsPPT}{聚苯基不对称三嗪模型化合物的质子化产物}
% \nomenclature{HMSPPT}{聚苯基对称三嗪模型化合物的质子化产物}
% \nomenclature{HMPPQ}{聚苯基喹噁啉模型化合物的质子化产物}
% \nomenclature{PDT}{热分解温度}
% \nomenclature{HPLC}{高效液相色谱 (High Performance Liquid Chromatography)}
% \nomenclature{HPCE}{高效毛细管电泳色谱 (High Performance Capillary lectrophoresis)}
% \nomenclature{LC-MS}{液相色谱-质谱联用 (Liquid chromatography-Mass Spectrum)}
% \nomenclature{TIC}{总离子浓度 (Total Ion Content)}
% \nomenclature{\textit{ab initio}}{基于第一原理的量子化学计算方法，常称从头算法}
% \nomenclature{DFT}{密度泛函理论 (Density Functional Theory)}
% \nomenclature{$E_a$}{化学反应的活化能 (Activation Energy)}
% \nomenclature{ZPE}{零点振动能 (Zero Vibration Energy)}
% \nomenclature{PES}{势能面 (Potential Energy Surface)}
% \nomenclature{TS}{过渡态 (Transition State)}
% \nomenclature{TST}{过渡态理论 (Transition State Theory)}
% \nomenclature{$\increment G^\neq$}{活化自由能(Activation Free Energy)}
% \nomenclature{$\kappa$}{传输系数 (Transmission Coefficient)}
% \nomenclature{IRC}{内禀反应坐标 (Intrinsic Reaction Coordinates)}
% \nomenclature{$\nu_i$}{虚频 (Imaginary Frequency)}
% \nomenclature{ONIOM}{分层算法 (Our own N-layered Integrated molecular Orbital and molecular Mechanics)}
% \nomenclature{SCF}{自洽场 (Self-Consistent Field)}
% \nomenclature{SCRF}{自洽反应场 (Self-Consistent Reaction Field)}



% 正文部分
\mainmatter
% !TeX root = ../main.tex

\chapter{学位论文形式结构}

\section{字数要求}
\begin{enumerate}
    \item 硕士论文: 正文一般为$1 \sim 3$万字；
    \item 博士论文: 正文一般不超过$15$万字。
\end{enumerate}


\section{前置部分}

\begin{enumerate}
    \item 封面和封底：由研究生院统一印刷，封面内容一律打印，其中导师姓名以研究生院备案名单为准。\\
          题目：在25字以内，能简明、具体、确切地表达论文特定内容，必要时，可以使用副标题。中文、英文题目应一致。
      \item 扉页内容包括学位论文中文、英文题目、专业名称、申请人姓名、导师姓名及论文答辩委员会组成 （由答辩委员会成员签名）。
      \item 原创性及学位论文使用授权声明。
      \item 中英文摘要。
\end{enumerate}


\section{主体部分}
\begin{enumerate}
    \item 主体部分包括引言 （前言），国内外文献综述，正文， 结语，参考文献。要求图表清晰，叙述流畅，章节有序，层次 分明。\\
    引言 （前言）部分内容主要为本研究课题的学术背景及理论与实际意义；本研究课题的来源及主要研究内容；建立研究的线索与思路。

    \item 文中的图、表、公式等，一律用阿拉伯数字按章顺序编号。如图 1-1、图2-2, 表 1-1、表 2-1，公式 (1-1) 等。图序及图名置于图的下方，居中排列；表序及表名置于表的上方，居中排列。详见第\ref{figures_tables}章的说明。

    \item 参考文献
    \begin{enumerate}
        \item 参考文献为论文中所有引文、引用观点以及对论文有重要影响和启发的文献；
        \item 参考文献按在论文中出现的先后依次排序；个别学科若通用该学科惯用的排序规范，可以例外；
        \item 参考文献内容一般排列在论文末尾 （论文篇幅较大且引用文献较多的，可在每章末尾注出），序码与论文加注处对应；
        \item 参考文献标注格式: 使用国标GB/T 7714-2015标准, 建议使用工具自动控制引文格式, 
        以保证格式规范。
        该\LaTeX{}模板已经对引文格式做了配置，
        用户需要将所需参考文献的信息存在 \texttt{bib} 格式的文件\texttt{ref/refs.bib}中, 
        通过 \verb|\cite{}|命令在恰当位置引用 (详见第\ref{citations}章的说明)。
        \texttt{bib}文件可从文献管理工具导出或自己用 \texttt{JabRef} 等软件编辑。
    \end{enumerate}

    \item 注释：可以用 “脚注”或 “文后注”来标注引用著作中的一些观点和案例，但全文标注方式应统一。
\end{enumerate}

\section{附录部分}
\begin{enumerate}
    \item 附录\\
        附录是正文主体的补充。下列内容可以作为附录：
        \begin{enumerate}
            \item 攻读学位期间发表的 （含已录用，并有录用通知书的）与学位论文相关的学术论文目录 。
            \item 由于篇幅过大，或取材于复制件不便编入正文的材料、数据。
            \item 对本专业同行有参考价值，但一般读者不必阅读的材料。
            \item 论文中使用的符号意义、单位缩写、程序全文及有关说明书。
            \item 附件：计算机程序清单、软磁盘、鉴定证书、获奖 奖状或专利证书的复印件等。
        \end{enumerate}
    \item 后记\\
        后记是有关本论文情况的说明性文字，主要是交代编写过程，阐述作者的感想和体会，对有关单位或个人的致谢语等。
\end{enumerate} 


% !TeX root = ../thuthesis-example.tex

\chapter{图表示例}

\section{插图}

图片通常在 \env{figure} 环境中使用 \cs{includegraphics} 插入，如图~\ref{fig:example} 的源代码。
建议矢量图片使用 PDF 格式，比如数据可视化的绘图；
照片应使用 JPG 格式；
其他的栅格图应使用无损的 PNG 格式。
注意，LaTeX 不支持 TIFF 格式；EPS 格式已经过时。

\begin{figure}
  \centering
  \includegraphics[width=0.6\linewidth]{example-image-a.pdf}
  \caption{示例图片}
  \label{fig:example}
\end{figure}

若图或表中有附注，采用英文小写字母顺序编号，附注写在图或表的下方。
% LaTeX 传统上一般将附注的内容同图表的标题写在一起，形成很长的一段文字。

如果一个图由两个或两个以上分图组成时，各分图分别以(a)、(b)、(c)...... 作为图序，并须有分图题。
推荐使用 \pkg{subcaption} 宏包来处理， 比如图~\ref{fig:subfig-a} 和图~\ref{fig:subfig-b}。

\begin{figure}
  \centering
  \subcaptionbox{分图 A\label{fig:subfig-a}}
    {\includegraphics[width=0.45\linewidth]{example-image-a.pdf}}
  \subcaptionbox{分图 B\label{fig:subfig-b}}
    {\includegraphics[width=0.45\linewidth]{example-image-b.pdf}}
  \caption{多个分图的示例}
  \label{fig:multi-image}
\end{figure}



\section{表格}

表应具有自明性。为使表格简洁易读，尽可能采用三线表，如表~\ref{tab:three-line}。
三条线可以使用 \pkg{booktabs} 宏包提供的命令生成。

\begin{table}
  \centering
  \caption{三线表示例}
  \begin{tabular}{ll}
    \toprule
    文件名          & 描述                         \\
    \midrule
    thuthesis.dtx   & 模板的源文件，包括文档和注释 \\
    thuthesis.cls   & 模板文件                     \\
    thuthesis-*.bst & BibTeX 参考文献表样式文件    \\
    thuthesis-*.bbx & BibLaTeX 参考文献表样式文件  \\
    thuthesis-*.cbx & BibLaTeX 引用样式文件        \\
    \bottomrule
  \end{tabular}
  \label{tab:three-line}
\end{table}

表格如果有附注，尤其是需要在表格中进行标注时，可以使用 \pkg{threeparttable} 宏包。
研究生要求使用英文小写字母 a、b、c……顺序编号，本科生使用圈码 ①、②、③……编号。

\begin{table}
  \centering
  \begin{threeparttable}[c]
    \caption{带附注的表格示例}
    \label{tab:three-part-table}
    \begin{tabular}{ll}
      \toprule
      文件名                 & 描述                         \\
      \midrule
      thuthesis.dtx\tnote{a} & 模板的源文件，包括文档和注释 \\
      thuthesis.cls\tnote{b} & 模板文件                     \\
      thuthesis-*.bst        & BibTeX 参考文献表样式文件    \\
      thuthesis-*.bbx        & BibLaTeX 参考文献表样式文件  \\
      thuthesis-*.cbx        & BibLaTeX 引用样式文件        \\
      \bottomrule
    \end{tabular}
    \begin{tablenotes}
      \item [a] 可以通过 xelatex 编译生成模板的使用说明文档；
        使用 xetex 编译 \file{thuthesis.ins} 时则会从 \file{.dtx} 中去除掉文档和注释，得到精简的 \file{.cls} 文件。
      \item [b] 更新模板时，一定要记得编译生成 \file{.cls} 文件，否则编译论文时载入的依然是旧版的模板。
    \end{tablenotes}
  \end{threeparttable}
\end{table}

\cleardoublepage

% !TeX root = ../main.tex

\chapter{引用文献的标注}\label{citations}

引用文献采用按照在文中被引用的顺序排序。
下文介绍 BibTeX 配合 \pkg{natbib} 宏包的主要使用方法。

在顺序编码制下，默认的 \cs{cite} 命令同 \cs{citep} 一样，序号置于方括号中，
引文页码会放在括号外。
统一处引用的连续序号会自动用短横线连接。

\sysusetup{
  cite-style = super,
}
\begin{tabular}{l@{\quad$\Rightarrow$\quad}l}
  \verb|\cite{zhangkun1994}|               & \cite{zhangkun1994}               \\
  \verb|\citet{zhangkun1994}|              & \citet{zhangkun1994}              \\
  \verb|\citep{zhangkun1994}|              & \citep{zhangkun1994}              \\
  \verb|\cite[42]{zhangkun1994}|           & \cite[42]{zhangkun1994}           \\
  \verb|\cite{zhangkun1994,zhukezhen1973}| & \cite{zhangkun1994,zhukezhen1973} \\
\end{tabular}


也可以取消上标格式，将数字序号作为文字的一部分。
建议全文统一使用相同的格式。

\begin{verbatim}
\sysusetup{
  cite-style = inline,
}
\end{verbatim}

\sysusetup{
  cite-style = inline,
}
\begin{tabular}{l@{\quad$\Rightarrow$\quad}l}
  \verb|\cite{zhangkun1994}|               & \cite{zhangkun1994}               \\
  \verb|\citet{zhangkun1994}|              & \citet{zhangkun1994}              \\
  \verb|\citep{zhangkun1994}|              & \citep{zhangkun1994}              \\
  \verb|\cite[42]{zhangkun1994}|           & \cite[42]{zhangkun1994}           \\
  \verb|\cite{zhangkun1994,zhukezhen1973}| & \cite{zhangkun1994,zhukezhen1973} \\
\end{tabular}



% \section{著者-出版年制}
% 
% 著者-出版年制下的 \cs{cite} 跟 \cs{citet} 一样。
% 
% \sysusetup{
%   cite-style = author-year,
% }
% \begin{tabular}{l@{\quad$\Rightarrow$\quad}l}
%   \verb|\cite{zhangkun1994}|                & \cite{zhangkun1994}                \\
%   \verb|\citet{zhangkun1994}|               & \citet{zhangkun1994}               \\
%   \verb|\citep{zhangkun1994}|               & \citep{zhangkun1994}               \\
%   \verb|\cite[42]{zhangkun1994}|            & \cite[42]{zhangkun1994}            \\
%   \verb|\citep{zhangkun1994,zhukezhen1973}| & \citep{zhangkun1994,zhukezhen1973} \\
% \end{tabular}

\vskip 2ex
\sysusetup{
  cite-style = super,
}
注意，引文参考文献的每条都要在正文中标注
\cite{zhangkun1994,zhukezhen1973,dupont1974bone,zhengkaiqing1987,%
  jiangxizhou1980,jianduju1994,merkt1995rotational,mellinger1996laser,%
  bixon1996dynamics,mahui1995,carlson1981two,taylor1983scanning,%
  taylor1981study,shimizu1983laser,atkinson1982experimental,%
  kusch1975perturbations,guangxi1993,huosini1989guwu,wangfuzhi1865songlun,%
  zhaoyaodong1998xinshidai,biaozhunhua2002tushu,chubanzhuanye2004,%
  who1970factors,peebles2001probability,baishunong1998zhiwu,%
  weinstein1974pathogenic,hanjiren1985lun,dizhi1936dizhi,%
  tushuguan1957tushuguanxue,aaas1883science,fugang2000fengsha,%
  xiaoyu2001chubanye,oclc2000about,scitor2000project%
}。

\cleardoublepage

% !TeX root = ../thuthesis-example.tex

\chapter{引用文献的标注}

模板支持 BibTeX 和 BibLaTeX 两种方式处理参考文献。
下文主要介绍 BibTeX 配合 \pkg{natbib} 宏包的主要使用方法。


\section{顺序编码制}

在顺序编码制下，默认的 \cs{cite} 命令同 \cs{citep} 一样，序号置于方括号中，
引文页码会放在括号外。
统一处引用的连续序号会自动用短横线连接。

\thusetup{
  cite-style = super,
}
\begin{tabular}{l@{\quad$\Rightarrow$\quad}l}
  \verb|\cite{zhangkun1994}|               & \cite{zhangkun1994}               \\
  \verb|\citet{zhangkun1994}|              & \citet{zhangkun1994}              \\
  \verb|\citep{zhangkun1994}|              & \citep{zhangkun1994}              \\
  \verb|\cite[42]{zhangkun1994}|           & \cite[42]{zhangkun1994}           \\
  \verb|\cite{zhangkun1994,zhukezhen1973}| & \cite{zhangkun1994,zhukezhen1973} \\
\end{tabular}


也可以取消上标格式，将数字序号作为文字的一部分。
建议全文统一使用相同的格式。

\thusetup{
  cite-style = inline,
}
\begin{tabular}{l@{\quad$\Rightarrow$\quad}l}
  \verb|\cite{zhangkun1994}|               & \cite{zhangkun1994}               \\
  \verb|\citet{zhangkun1994}|              & \citet{zhangkun1994}              \\
  \verb|\citep{zhangkun1994}|              & \citep{zhangkun1994}              \\
  \verb|\cite[42]{zhangkun1994}|           & \cite[42]{zhangkun1994}           \\
  \verb|\cite{zhangkun1994,zhukezhen1973}| & \cite{zhangkun1994,zhukezhen1973} \\
\end{tabular}



\section{著者-出版年制}

著者-出版年制下的 \cs{cite} 跟 \cs{citet} 一样。

\thusetup{
  cite-style = author-year,
}
\begin{tabular}{l@{\quad$\Rightarrow$\quad}l}
  \verb|\cite{zhangkun1994}|                & \cite{zhangkun1994}                \\
  \verb|\citet{zhangkun1994}|               & \citet{zhangkun1994}               \\
  \verb|\citep{zhangkun1994}|               & \citep{zhangkun1994}               \\
  \verb|\cite[42]{zhangkun1994}|            & \cite[42]{zhangkun1994}            \\
  \verb|\citep{zhangkun1994,zhukezhen1973}| & \citep{zhangkun1994,zhukezhen1973} \\
\end{tabular}

\vskip 2ex
\thusetup{
  cite-style = super,
}
注意，引文参考文献的每条都要在正文中标注
\cite{zhangkun1994,zhukezhen1973,dupont1974bone,zhengkaiqing1987,%
  jiangxizhou1980,jianduju1994,merkt1995rotational,mellinger1996laser,%
  bixon1996dynamics,mahui1995,carlson1981two,taylor1983scanning,%
  taylor1981study,shimizu1983laser,atkinson1982experimental,%
  kusch1975perturbations,guangxi1993,huosini1989guwu,wangfuzhi1865songlun,%
  zhaoyaodong1998xinshidai,biaozhunhua2002tushu,chubanzhuanye2004,%
  who1970factors,peebles2001probability,baishunong1998zhiwu,%
  weinstein1974pathogenic,hanjiren1985lun,dizhi1936dizhi,%
  tushuguan1957tushuguanxue,aaas1883science,fugang2000fengsha,%
  xiaoyu2001chubanye,oclc2000about,scitor2000project%
}。



% 其他部分
\backmatter

% 参考文献
\bibliography{ref/refs}  % 参考文献使用 BibTeX 编译
% \printbibliography       % 参考文献使用 BibLaTeX 编译

% 附录
\appendix
\chapter{补充内容}

附录是与论文内容密切相关、但编入正文又影响整篇论文编排的条理和逻辑性的资料，例如某些重要的数据表格、计算程序、统计表等，是论文主体的补充内容，可根据需要设置。


\section{图表示例}

\subsection{图}

附录中的图片示例（图~\ref{fig:appendix-figure}）。

\begin{figure}
  \centering
  \includegraphics[width=0.6\linewidth]{example-image-a.pdf}
  \caption{附录中的图片示例}
  \label{fig:appendix-figure}
\end{figure}


\subsection{表格}

附录中的表格示例（表~\ref{tab:appendix-table}）。

\begin{table}
  \centering
  \caption{附录中的表格示例}
  \begin{tabular}{ll}
    \toprule
    文件名          & 描述                         \\
    \midrule
    thuthesis.dtx   & 模板的源文件，包括文档和注释 \\
    thuthesis.cls   & 模板文件                     \\
    thuthesis-*.bst & BibTeX 参考文献表样式文件    \\
    thuthesis-*.bbx & BibLaTeX 参考文献表样式文件  \\
    thuthesis-*.cbx & BibLaTeX 引用样式文件        \\
    \bottomrule
  \end{tabular}
  \label{tab:appendix-table}
\end{table}


\section{数学公式}

附录中的数学公式示例（公式~\eqref{eq:appendix-equation}）。
\begin{equation}
  \frac{1}{2 \symup{\pi} \symup{i}} \int_\gamma f = \sum_{k=1}^m n(\gamma; a_k) \mathscr{R}(f; a_k)
  \label{eq:appendix-equation}
\end{equation}

% % !TeX root = ../thuthesis-example.tex

\begin{survey}
\label{cha:survey}

\title{Title of the Survey}
\maketitle


\tableofcontents


本科生的外文资料调研阅读报告。


\section{Figures and Tables}

\subsection{Figures}

An example figure in appendix (Figure~\ref{fig:appendix-survey-figure}).

\begin{figure}
  \centering
  \includegraphics[width=0.6\linewidth]{example-image-a.pdf}
  \caption{Example figure in appendix}
  \label{fig:appendix-survey-figure}
\end{figure}


\subsection{Tables}

An example table in appendix (Table~\ref{tab:appendix-survey-table}).

\begin{table}
  \centering
  \caption{Example table in appendix}
  \begin{tabular}{ll}
    \toprule
    File name       & Description                                         \\
    \midrule
    thuthesis.dtx   & The source file including documentaion and comments \\
    thuthesis.cls   & The template file                                   \\
    thuthesis-*.bst & BibTeX styles                                       \\
    thuthesis-*.bbx & BibLaTeX styles for bibliographies                  \\
    thuthesis-*.cbx & BibLaTeX styles for citations                       \\
    \bottomrule
  \end{tabular}
  \label{tab:appendix-survey-table}
\end{table}


\section{Equations}

An example equation in appendix (Equation~\eqref{eq:appendix-survey-equation}).
\begin{equation}
  \frac{1}{2 \symup{\pi} \symup{i}} \int_\gamma f = \sum_{k=1}^m n(\gamma; a_k) \mathscr{R}(f; a_k)
  \label{eq:appendix-survey-equation}
\end{equation}


\section{Citations}

Example citations in appendix.
\cite{abrahams99tex}
\cite{salomon1995advanced}
\cite{abrahams99tex,salomon1995advanced}


\bibliographystyle{unsrtnat}
\bibliography{ref/appendix}

\end{survey}
       % 本科生：外文资料的调研阅读报告
% % !TeX root = ../thuthesis-example.tex

\begin{translation}
\label{cha:translation}

\title{书面翻译题目}
\maketitle

\tableofcontents


本科生的外文资料书面翻译。


\section{图表示例}

\subsection{图}

附录中的图片示例（图~\ref{fig:appendix-translation-figure}）。

\begin{figure}
  \centering
  \includegraphics[width=0.6\linewidth]{example-image-a.pdf}
  \caption{附录中的图片示例}
  \label{fig:appendix-translation-figure}
\end{figure}


\subsection{表格}

附录中的表格示例（表~\ref{tab:appendix-translation-table}）。

\begin{table}
  \centering
  \caption{附录中的表格示例}
  \begin{tabular}{ll}
    \toprule
    文件名          & 描述                         \\
    \midrule
    thuthesis.dtx   & 模板的源文件，包括文档和注释 \\
    thuthesis.cls   & 模板文件                     \\
    thuthesis-*.bst & BibTeX 参考文献表样式文件    \\
    thuthesis-*.bbx & BibLaTeX 参考文献表样式文件  \\
    thuthesis-*.cbx & BibLaTeX 引用样式文件        \\
    \bottomrule
  \end{tabular}
  \label{tab:appendix-translation-table}
\end{table}


\section{数学公式}

附录中的数学公式示例（公式~\eqref{eq:appendix-translation-equation}）。
\begin{equation}
  \frac{1}{2 \symup{\pi} \symup{i}} \int_\gamma f = \sum_{k=1}^m n(\gamma; a_k) \mathscr{R}(f; a_k)
  \label{eq:appendix-translation-equation}
\end{equation}


\section{文献引用}

文献引用示例\cite{abrahams99tex}。


% 书面翻译的参考文献
\bibliographystyle{unsrtnat}
\bibliography{ref/appendix}

% 书面翻译对应的原文索引
\begin{translation-index}
  \nocite{salomon1995advanced}
  \bibliographystyle{unsrtnat}
  \bibliography{ref/appendix}
\end{translation-index}

\end{translation}
  % 本科生：外文资料的书面翻译

% 致谢
% !TeX root = ../thuthesis-example.tex

\begin{acknowledgements}
  衷心感谢导师×××教授和物理系××副教授对本人的精心指导。他们的言传身教将使我终生受益。

  在美国麻省理工学院化学系进行九个月的合作研究期间，承蒙 Robert Field 教授热心指导与帮助，不胜感激。

  感谢×××××实验室主任×××教授，以及实验室全体老师和同窗们学的热情帮助和支持！

  本课题承蒙国家自然科学基金资助，特此致谢。
\end{acknowledgements}



% 个人简历、在学期间完成的相关学术成果
% !TeX root = ../thuthesis-example.tex

\begin{resume}

  \section*{个人简历}

  197× 年 ×× 月 ×× 日出生于四川××县。

  1992 年 9 月考入××大学化学系××化学专业，1996 年 7 月本科毕业并获得理学学士学位。

  1996 年 9 月免试进入清华大学化学系攻读××化学博士至今。


  \section*{在学期间完成的相关学术成果}

  \subsection*{学术论文}

  \begin{achievements}
    \item Yang Y, Ren T L, Zhang L T, et al. Miniature microphone with silicon-based ferroelectric thin films[J]. Integrated Ferroelectrics, 2003, 52:229-235.
    \item 杨轶, 张宁欣, 任天令, 等. 硅基铁电微声学器件中薄膜残余应力的研究[J]. 中国机械工程, 2005, 16(14):1289-1291.
    \item 杨轶, 张宁欣, 任天令, 等. 集成铁电器件中的关键工艺研究[J]. 仪器仪表学报, 2003, 24(S4):192-193.
    \item Yang Y, Ren T L, Zhu Y P, et al. PMUTs for handwriting recognition. In press[J]. (已被Integrated Ferroelectrics录用)
  \end{achievements}


  \subsection*{专利}

  \begin{achievements}
    \item 任天令, 杨轶, 朱一平, 等. 硅基铁电微声学传感器畴极化区域控制和电极连接的方法: 中国, CN1602118A[P]. 2005-03-30.
    \item Ren T L, Yang Y, Zhu Y P, et al. Piezoelectric micro acoustic sensor based on ferroelectric materials: USA, No.11/215, 102[P]. (美国发明专利申请号.)
  \end{achievements}

\end{resume}


% % 指导教师/指导小组学术评语
% % !TeX root = ../thuthesis-example.tex

\chapter{指导小组学术评语}

论文提出了……


% % 答辩委员会决议书
% % !TeX root = ../thuthesis-example.tex

\chapter{答辩委员会决议书}

论文提出了……

论文取得的主要创新性成果包括：

1. ……

2. ……

3. ……

论文工作表明作者在×××××具有×××××知识，具有××××能力，论文××××，答辩××××。

答辩委员会表决，（×票/一致）同意通过论文答辩，并建议授予×××（姓名）×××（门类）学博士/硕士学位。


% 本科生的综合论文训练记录表（扫描版）
% \record{file=scan-record.pdf}

\end{document}
